\documentclass[11pt,a4paper]{article}

\usepackage[utf8]{inputenc}
\usepackage[T1]{fontenc}
\usepackage{graphicx}
\usepackage{listings}
\usepackage{xcolor}
\usepackage{booktabs}
\usepackage{geometry}
\usepackage{enumitem}
\geometry{margin=2.5cm}

% Le pack de langue 'french' pour babel n'est pas installe dans cet
% environnement de test minimal ; sur Overleaf / TeXLive complet,
% decommenter la ligne suivante pour les regles typographiques francaises :
% \usepackage[french]{babel}

\lstset{basicstyle=\ttfamily\small, breaklines=true, frame=single}

\title{Manuel utilisateur \\ \large Catalogue d'un commerce en ligne}
\author{[Nom 1] \and [Nom 2] \\ Licence 2 Mathématiques-Informatique \\ UIDT Thiès}
\date{Année universitaire 2025--2026}

\begin{document}
\maketitle

\section{Présentation}

Ce programme en C implémente un catalogue de commerce en ligne (produits,
clients, commandes) sur trois structures de données distinctes (tableau
statique, tableau dynamique de pointeurs, liste doublement chaînée), avec
persistance automatique entre deux exécutions.

\section{Compilation}

\begin{lstlisting}
make            # compile le programme (cible "all")
make benchmark  # compile en -O2 et lance directement les mesures
make clean      # supprime l'executable et les fichiers de donnees .dat
\end{lstlisting}

Compilation manuelle équivalente :
\begin{lstlisting}
gcc -Wall -Wextra -std=c11 src/PROJETSD1.c -o projetsd1 -lm
\end{lstlisting}

Le fichier \texttt{produits.h} doit se trouver dans le même dossier que
\texttt{PROJETSD1.c}.

\section{Lancement}

\begin{lstlisting}
./projetsd1
\end{lstlisting}

Au premier lancement, aucun fichier de données n'existe encore : le
programme l'indique (\textit{"Aucun fichier 'produits.dat' trouvé. Catalogue
vide."}) et démarre avec un catalogue vide. Aux lancements suivants, toutes
les données précédemment enregistrées sont rechargées automatiquement.

\section{Menu principal}

\begin{table}[h]
\centering
\begin{tabular}{cl}
\toprule
Option & Fonction \\
\midrule
1  & Gestion du catalogue (ajouter / afficher un produit) \\
2  & Créer un compte client / afficher la liste des clients \\
3  & Ajouter un produit au panier / voir le catalogue \\
4  & Voir son panier / passer une commande \\
5  & Voir toutes les commandes passées \\
6  & Statistiques (min, max, moyenne, médiane, écart-type des prix) \\
7  & Trier et afficher le tableau statique (par prix croissant) \\
8  & Trier et afficher le tableau dynamique (par stock décroissant) \\
9  & Gestion de la liste doublement chaînée (sous-menu détaillé, §6) \\
10 & Mise à jour d'un produit existant \\
0  & Quitter (sauvegarde automatique avant fermeture) \\
\bottomrule
\end{tabular}
\caption{Options du menu principal}
\end{table}

\section{Gestion du catalogue (option 1)}

\begin{enumerate}
  \item \textbf{Ajouter un produit} : demande successivement l'identifiant
        (entier unique, \textbf{ne jamais utiliser 0}, réservé en interne
        pour repérer la fin du tableau statique), le nom, la catégorie, le
        prix et le stock. Le produit est ajouté simultanément au tableau
        statique \emph{et} au tableau dynamique.
  \item \textbf{Afficher le catalogue} : liste tous les produits du tableau
        statique sous forme de tableau formaté (ID, nom, catégorie, prix, stock).
\end{enumerate}

\section{Liste doublement chaînée (option 9)}

Sous-menu dédié, indépendant du tableau statique/dynamique (un produit doit
être saisi séparément pour apparaître dans la liste) :

\begin{table}[h]
\centering
\begin{tabular}{cl}
\toprule
Sous-option & Fonction \\
\midrule
1 & Ajouter un produit en tête de liste \\
2 & Ajouter un produit en queue de liste \\
3 & Afficher la liste \\
4 & Rechercher un produit par identifiant \\
5 & Recherche multicritère (catégorie exacte + intervalle de prix) \\
6 & Supprimer un produit par identifiant \\
7 & Trier la liste par prix croissant (tri par insertion) puis l'afficher \\
8 & Statistiques sur les prix de la liste \\
\bottomrule
\end{tabular}
\caption{Sous-menu "Liste doublement chaînée"}
\end{table}

\section{Clients, panier et commandes (options 2 à 5)}

\begin{itemize}
  \item \textbf{Option 2} : créer un client (identifiant, nom, adresse) ou
        afficher la liste des clients enregistrés.
  \item \textbf{Option 3} : ajouter un produit du catalogue au panier d'un
        client (par identifiant client puis identifiant produit), ou
        consulter le catalogue complet avant de choisir.
  \item \textbf{Option 4} : afficher le panier en cours d'un client
        (identifiant requis), ou transformer ce panier en commande
        enregistrée (un identifiant de commande est attribué automatiquement).
  \item \textbf{Option 5} : afficher l'historique de toutes les commandes
        passées, avec le client et les produits associés.
\end{itemize}

\section{Statistiques (option 6)}

Calcule, sur le tableau statique \emph{ou} le tableau dynamique (au choix),
les agrégations suivantes sur les prix : minimum, maximum, moyenne,
médiane et écart-type.

\section{Mise à jour d'un produit (option 10)}

Demande l'identifiant produit puis, selon la structure choisie (statique,
dynamique ou liste), permet de modifier un seul champ à la fois : nom,
catégorie, prix ou stock.

\section{Persistance des données}

Toutes les données sont sauvegardées automatiquement à la fermeture
(option 0) dans des fichiers binaires situés dans le dossier d'exécution :

\begin{table}[h]
\centering
\begin{tabular}{ll}
\toprule
Fichier & Contenu \\
\midrule
\texttt{produits.dat}   & Tableau statique des produits \\
\texttt{clients.dat}    & Liste des clients \\
\texttt{commandes.dat}  & Historique des commandes \\
\texttt{liste.dat}      & Contenu de la liste doublement chaînée \\
\texttt{dynamique.dat}  & Contenu du tableau dynamique de pointeurs \\
\bottomrule
\end{tabular}
\caption{Fichiers de persistance générés}
\end{table}

Pour repartir d'un catalogue vide, supprimez ces fichiers (ou lancez
\texttt{make clean}).

\section{Mode benchmark}

Une commande dédiée mesure les performances des trois structures :

\begin{lstlisting}
./projetsd1 --benchmark
\end{lstlisting}

Ce mode génère des jeux de données de tailles croissantes (jusqu'à
$n=10^5$), selon trois distributions (aléatoire, triée, inversée), répète
chaque mesure 10 fois, et écrit :
\begin{itemize}
  \item \texttt{resultats\_benchmark.csv} : les temps mesurés (insertion,
        recherche, tri) pour chaque combinaison ;
  \item \texttt{data/jeu\_n<taille>\_<distribution>.csv} : un échantillon des
        données générées (pour $n \le 1000$).
\end{itemize}

Les courbes correspondantes s'obtiennent ensuite avec :
\begin{lstlisting}
python3 scripts/generer_courbes.py resultats_benchmark.csv
\end{lstlisting}

\textbf{Attention} : à $n=10^5$, le tri par insertion (coût $O(n^2)$ dans le
pire cas) peut prendre plusieurs minutes au total ; c'est attendu.

\section{Dépannage}

\begin{itemize}
  \item \textit{"Produit introuvable"} : vérifier que l'identifiant existe
        bien dans la structure interrogée (un produit ajouté uniquement via
        le sous-menu liste, option 9, n'existe pas dans le tableau statique,
        et inversement).
  \item Saisie d'un texte au lieu d'un nombre attendu : le programme affiche
        \textit{"Saisie invalide"} et redemande une valeur au menu principal ;
        dans les sous-menus, une saisie non numérique peut nécessiter de
        relancer le programme.
  \item Pour repartir d'un état totalement vierge : supprimer tous les
        fichiers \texttt{*.dat} avant de relancer.
\end{itemize}

\end{document}
